\documentclass[11pt,a4paper]{article}
\usepackage[utf8]{inputenc}
\usepackage[french]{babel}
\usepackage[T1]{fontenc}
\usepackage{lmodern}
\usepackage[margin=1.5cm]{geometry}

\begin{document}

% En-tête
\begin{center}
    {\Huge \textbf{Samy Dubois}} \\[0.2cm]
    {\Large \textbf{Ingénieur DevOps}} \\[0.1cm]
    samy.dubois@email.com \ | \ 06 72 56 83 53
\end{center}

\vspace{0.3cm}

% Résumé / Profil
\section*{Profil Professionnel}
\noindent
Site Reliability Engineer passionné par la haute disponibilité et la performance des systèmes distribués. J'interviens sur la sécurisation, le monitoring et la scalabilité des serveurs cloud pour garantir un taux de disponibilité (SLA) de 99,99\% sur les applications critiques. Mon but est de construire des infrastructures capables de résister à n'importe quel incident.

\vspace{0.3cm}

% Compétences
\section*{Compétences Techniques}
\noindent
\textbf{Python} $\bullet$ \textbf{GitLab CI/CD} $\bullet$ \textbf{Ansible} $\bullet$ \textbf{AWS} $\bullet$ \textbf{Bash} $\bullet$ \textbf{Prometheus} $\bullet$ \textbf{Linux} $\bullet$ \textbf{Kubernetes}

\vspace{0.3cm}

% Expériences
\section*{Expériences Professionnelles}
\noindent \textbf{Ingénieur DevOps / Confirmé} --- \textit{Orange Business Services} \hfill \textbf{12/2022 à Aujourd'hui} \\
Astreinte technique (On-call) et résolution d'incidents critiques en production de niveau 3 (N3), avec rédaction de post-mortems détaillés (blameless). Gestion et sécurisation des secrets de l'entreprise via HashiCorp Vault, avec rotation automatique des clés d'accès des API tierces. Mise en place d'une architecture de journalisation centralisée (ELK Stack), facilitant l'analyse des logs et le diagnostic des anomalies en production. Conception d'une infrastructure cloud résiliente sur AWS via Terraform (IaC), assurant une haute disponibilité des services et une scalabilité horizontale automatique. \\[0.3cm]
\noindent \textbf{Ingénieur DevOps} --- \textit{OVHcloud} \hfill \textbf{08/2020 à 12/2022} \\
Mise en place de pipelines d'intégration et de déploiement continus (CI/CD) sur GitLab automatisant l'ensemble des tests et le déploiement sur les clusters Kubernetes. Déploiement d'une politique de sécurité 'Zero Trust', incluant la gestion stricte des identités (IAM) et le chiffrement des données au repos. Audit de l'infrastructure cloud et application des principes du FinOps, permettant une réduction de la facture mensuelle AWS de 25\%. Automatisation du provisionnement et de la configuration des serveurs Linux à l'aide de playbooks Ansible, éliminant les erreurs manuelles. \\[0.3cm]
\noindent \textbf{Alternance — Assistant Ingénieur DevOps} --- \textit{Datadog} \hfill \textbf{04/2018 à 04/2020} \\
Migration de 50 applications legacy vers une architecture conteneurisée sous Docker, réduisant drastiquement les coûts d'hébergement annuels. Implémentation d'un système de monitoring complet avec Prometheus et Grafana, réduisant le temps de détection des pannes (MTTD) de 60\%. Mise en place d'une stratégie de sauvegarde automatisée (Disaster Recovery Plan) avec des tests de restauration trimestriels réussis à 100\%. \\[0.3cm]


\vspace{0.3cm}

% Formations
\section*{Formation \& Diplômes}
\noindent \textbf{Master en Cloud Computing} \hfill \textbf{2018 - 2020} \\
\textit{Sorbonne Université} \\[0.3cm]
\noindent \textbf{Licence en Cloud Computing} \hfill \textbf{2015 - 2018} \\
\textit{Sorbonne Université} \\[0.3cm]


\end{document}